% PRJ93 literature review — restructured draft, revision 3 (2026-08-02).
% Working copy under brain/. NOT pushed to Overleaf (gate 5).
%
% Revision 3 responds to critique iteration 2 (Roles A/B/C + synthesis).
% The full record, including what was rejected and why, is in
% brain/ledger/litreview_critique.md.
%
% Drafting rule adopted at revision 3, after the iteration-2 synthesis
% identified the chapter's defect generator: this chapter states what the
% literature says and what this project's design assumes. It makes no
% statement about what any later chapter will do, and no claim about the
% author's own rigour that a reader cannot check. Where a fact is known it is
% stated; where it is not known, nothing is promised in its place.
%
% No project result figure appears in this chapter.

This chapter is organised as a methodological pipeline that mirrors the
proposed system: an agent must first learn a venue's normal trading rhythm,
then satisfy itself that the rhythm has been scored and selected honestly, then
notice when the present departs from it, then decide whether the departure is
worth raising, and finally be judged on whether that decision was a good one.
Each section consumes the output of the one before it. Two of them are
deliberately parallel: Section~\ref{sec:rw-ruler} asks how a forecaster should
be judged, and Section~\ref{sec:rw-evaluation} asks the same question of the
agent built on top of it.

Four arguments are threaded through the whole, and this review takes a position
on each rather than surveying opinion. Operational rhythm can be borrowed
across venues, but the established theory of that benefit concerns the fitting
of shared models over large collections, and its extension to a three-venue
estate is this chapter's own conjecture rather than an inherited result. A
deviation is best defined as the breach of a calibrated band, whose guarantee
degrades gracefully rather than holding through a regime change. A learned
rhythm functions as the agent's model of normality, in a bounded sense argued
for below rather than assumed. And the worth of a proactive alert turns on the
asymmetric human cost of a false alarm, which is an empirical claim with
empirical support and is treated here as one.

Three features of the evidence base are stated here and recalled where they
bear. First, the frontier of proactive agency is young. Twelve of the works
cited in this chapter had not completed peer review at the time of writing:
eleven arXiv preprints dated between 2024 and 2026, and one conference
submission. One is in Section~\ref{sec:rw-framing}, three in
Section~\ref{sec:rw-rhythm} --- including the foundation model this work
ultimately serves --- and eight across Sections~\ref{sec:rw-surfacing}
and~\ref{sec:rw-evaluation}. Each of the twelve is marked at the point where it
carries an argument; other recent work is cited without that mark where it is
corroborative rather than load-bearing. Second, the contribution claim
developed in Section~\ref{sec:rw-synthesis} is positioned against unrefereed
work: the system it distinguishes itself from, and the metric template it
adapts, are both 2026 preprints. That is a stated risk of working at this
frontier, and Section~\ref{sec:rw-synthesis} says what changes if those results
do not survive review. Third, this review was written after the experimental
work it introduces. Where the literature reviewed here anticipates an outcome
in this estate, that anticipation is offered as an argument the reader can
check against the source, not as a record of what was believed beforehand. Two
places where the literature points the other way from what this project
observed are marked as such, in Sections~\ref{sec:rw-ruler}
and~\ref{sec:rw-deviation}.

\section{From Decision Support to the Autonomous Managerial Agent}
\label{sec:rw-framing}

The system this dissertation builds is a manager-facing agent, and a
related-work chapter must first settle what kind of artefact that makes it.
\citet{gorry_framework_1971} drew the distinction that
still organises the field: a decision is \emph{structured} when Simon's
intelligence, design and choice phases can be formalised into rules, and
\emph{unstructured} when it resists formalisation and demands human judgement,
with decision-support systems positioned to assist managers on the
unstructured and semi-structured cases rather than to replace them. For five
decades the manager remained the decision-maker and the software remained the
support. \citet{kumar_agentic_2026} argue that agentic AI breaks
this arrangement. What separates an agent from earlier enterprise software,
on their account, is delegated autonomy: the capacity to initiate actions,
make decisions and coordinate activity toward a goal with little human
prompting, so that agency is distributed across people and machines rather
than held by the user alone. Read against Gorry, this is not a faster support
tool but a different category of actor, one that participates in the decision
instead of merely informing it.

The deployment record indicates the category is already populated, and that it
is populated faster than it is understood. An index of thirty widely used
agentic products \citep[a 2026 preprint]{staufer_2025_2026} documents systems
acting with limited oversight, yet also finds that developers disclose almost
nothing about their safety and evaluation, leaving no public information for
135 of 240 safety fields. The gap is not peculiar to agents. Surveying
published reports of machine-learning deployments across industries and use
cases, \citet{paleyes_challenges_2022} map the difficulties practitioners
report onto the stages of a deployment workflow and find that practitioners
face issues at every stage of that process --- a catalogue assembled from case
studies of systems in use, not from benchmark results. Two consequences follow
for this work. Principled evaluation of a deployed agent is a contribution in
its own right rather than an afterthought, and the setting in which such an
agent is evaluated is itself a variable, because the difficulties a deployment
reports are drawn from a different sample than the difficulties a benchmark
reports.

This dissertation takes the position, which the remainder of the chapter
develops rather than assumes, that an agent acting unasked must hold an
internal model of what is \emph{normal} for the venue it manages, because an
unprompted action is warranted only by a departure from expectation. The
premise is not idle: the evidence for it arrives in
Section~\ref{sec:rw-surfacing}, where the dominant documented failure of
proactive systems is over-offering \citep{lu_proactive_2024}, and in
Section~\ref{sec:rw-evaluation}, where the cost of that failure is shown to
fall on the operator's willingness to keep reading
\citep{dixon_independence_2007,ancker_effects_2017}. Learning a model of
normality from short, noisy trading records is the constraint, named here and
resolved later, that shadows the whole design.

\section{Learning the Rhythm: Forecasting Normality Under Data Scarcity}
\label{sec:rw-rhythm}

By a venue's \emph{rhythm} we mean its recurring pattern of trade, the
day-of-week and seasonal regularities against which any given day can be
called typical or surprising. The defining difficulty is scarcity: a single
hospitality venue offers months, not years, of clean history, and a venue
that has just opened offers none at all. This is the cold-start constraint,
and it rules out the comfortable assumption of much forecasting work that each
series carries enough of its own history to be fitted in isolation.

Within that constraint, the demand-forecasting literature delivers a result
that is easy to overlook: simple methods are hard to beat.
\citet{chae_value_2024}, studying a large restaurant chain, show that deep
models using only internal operational and calendar features match or
outperform machine-learning models given access to richer external
macroeconomic and health data, so the marginal value of additional data is
smaller than practitioners assume. \citet{hossain_comparative_2025}
reinforce the point from the opposite direction: although a multilayer
perceptron and random forest rank highest on their restaurant data, classical
exponential smoothing and \citet{croston_forecasting_1972}'s method outperform
a gradient-boosted baseline, which fails to capture seasonality without
hand-built features. \citet{schmidt_machine_2022} report a related pattern on a
single restaurant's sales: across twenty-three models and two baselines, the
best one-day-ahead result is a kernel-ridge model at 19.6\% sMAPE, while at the
one-week horizon the non-recurrent models degrade past 20\% error and the best
recurrent model reports 19.5\%. The separations at one week are a fraction of a
percentage point on one venue with no dispersion reported, so this is an
illustration of a possible horizon dependence and not evidence of one --- a
caution the rank-stability literature of Section~\ref{sec:rw-ruler} makes
general. Taken together these studies establish a baseline expectation this
work inherits: on small, seasonal hospitality series, classical statistical and
linear models set a competitive bar that a more elaborate method must clear on
evidence rather than on architecture.

\subsection*{When does borrowing across series pay?}

The transfer route promises to relieve scarcity by importing structure learned
elsewhere, and the question this section must settle is not whether that is
possible but when it is worth doing.
\citet{montero-manso_principles_2021} give the clearest available answer for
the case they treat, and it is not the one the enthusiasm for shared models
suggests. Globality is not a similarity assumption: global methods are, in
their words, not more restrictive than local methods, since both can produce
the same forecasts without any assumptions about similarity of the series.
What differs is where the estimation budget is spent. The complexity of local
methods grows with the size of the set while it remains constant for global
methods, so that --- and the qualifier is the authors' own --- \emph{in large
datasets} a global algorithm can afford to be quite complex and still benefit
from better generalisation. The benefit of pooling, on this account, is a
statement about the number of series available to amortise a shared function
over.

Two limits on that result must be stated before it is used here. It is a claim
about \emph{estimation} --- the complexity budget of fitting a shared function
across a collection --- whereas the mechanism this work ultimately employs
conditions a frozen pretrained model on sibling series at inference time,
fitting nothing on the estate at all. And the formalism supplies no threshold:
nothing in it says that three series is small and forty is large. The
expectation this dissertation carries into its experiments --- that pooling
across a three-venue estate will yield little --- is therefore an extension of
the principle by analogy, and is this chapter's conjecture rather than
\citeauthor{montero-manso_principles_2021}'s result. It is stated so that it
can fail: the analogy is wrong if in-context conditioning is found to help at
group sizes where shared estimation would not, and a positive transfer result
on a target collection of a handful of series would be evidence against it. The
three demonstrations relied on below are not such cases --- in-context
fine-tuning is evaluated on established multi-series forecasting benchmarks
\citep{das_-context_2025}, cold-start web-traffic retrieval draws on a
collection of pages \citep{zhou_context-driven_2025}, and GPHT pretrains on a
mixed corpus assembled from several public datasets
\citep{liu_generative_2024} --- but no systematic search for a small-collection
counterexample was undertaken, and none is claimed.

With that expectation set, the mechanisms are worth stating precisely.
Time-series foundation models pretrain on large heterogeneous corpora and then
forecast unseen series zero-shot, and they divide by representation. TimesFM
\citep{das_decoder-only_2024} is a patched decoder that predicts the next patch
autoregressively and reaches near-supervised zero-shot accuracy with 200M
parameters. Chronos \citep{ansari_chronos_2024} takes the opposite stance,
scaling and quantising real values into a fixed token vocabulary so an
off-the-shelf language-model architecture can be trained on them; on the 27
held-out datasets of its 42-dataset benchmark it reports zero-shot accuracy,
measured by weighted quantile loss and by mean absolute scaled error,
comparable or superior to models trained on those series directly. A wider
family --- masked-encoder, quantile, masked-reconstruction and probabilistic or
industrial entrants
\citep{woo_unified_2024,liu_moirai_2026,goswami_moment_2024,rasul_lag-llama_2024,garza_timegpt-1_2024}
--- fills out the design space without changing the central premise that
pretraining on time series, not on language, is what transfers.

That premise needs a sharp qualification, and \citet{tan_are_2024} supply it.
In ablations on three popular methods that bolt a pretrained language model
onto a forecaster, they find that removing the language model or replacing it
with a basic attention layer leaves accuracy unchanged or better while cutting
training and inference time by up to three orders of magnitude, and that the
language model helps neither sequence modelling nor the few-shot regime. This
does not contradict TimesFM and Chronos so much as locate the value precisely:
gains come from time-series pretraining and patching, not from language-model
machinery. TimesFM reaches the same conclusion from the other direction,
reporting an improvement of more than 25\% on its scaled error metric over a
language-model prompting baseline \citep{das_decoder-only_2024}. For a small
hospitality estate the lesson is a disciplined baseline ladder, classical
models first, then a transfer-based global model, with scepticism toward any
component justified by language-model provenance alone.

One entrant deserves separate treatment, because it is the model this work
eventually serves. Chronos-2 \citep[a 2025 preprint]{ansari_chronos-2_2025}
departs from its predecessor by alternating standard time attention with a
\emph{group} attention layer that aggregates across all series sharing a group
identifier at each patch index. A group may be a set of series sharing a source
or metadata, so that the model performs cross-series learning at inference time
instead of forecasting each series from its own history alone, and the authors
single out short histories and cold starts as the conditions under which that
sharing helps most. Read against the conjecture above, that claim is one this
dissertation tests rather than inherits: group attention supplies a mechanism
for pooling, and whether a group of three is large enough for the mechanism to
pay is precisely the open question.

The same mechanism admits past-only and known-future covariates as group
members, which places Chronos-2 among the small number of pretrained
forecasters that can condition on exogenous information at all. TabPFN-TS
\citep[a 2026 preprint]{hoo_tables_2026}, an 11M-parameter tabular foundation
model applied to forecasting through temporal featurisation, is the other, and
it inherits from TabPFN \citep{hollmann_accurate_2025} a design explicitly
aimed at datasets of up to ten thousand rows. \citet{hertel_explainable_2026}
evaluate both on load forecasting and find them competitive with transformers
trained from scratch on domain-specific data while requiring no training data
at all. A single hospitality estate, with three venues, one year of trade and a
handful of calendar and weather regressors, sits inside the regime these two
models were designed for rather than at the scale the transformer literature
usually assumes.

This dissertation assumes, rather than demonstrates, that the covariates most
worth conditioning on in hospitality are known-future calendar variables ---
the day of week, the bank holiday, the fixture in the diary --- rather than
weather. No source surveyed here separates a calendar channel from a weather
channel in a way that would establish the split; the assumption is recorded so
that it can be disputed. It matters because the next subsection argues that the
weather covariate specifically should be expected to contribute little, and a
forecaster chosen for covariate capability alone would then be chosen on a
capability the following argument discounts.

What makes transfer usable under genuine cold-start is the ability to adapt at
inference without per-venue retraining. \citet{das_-context_2025} demonstrate
in-context fine-tuning, prompting a foundation model with several related
series so it adapts to a target distribution at inference time, and report a
best-case improvement of up to 25\% on their scaled error metric over competing
baselines while rivalling a model explicitly fine-tuned on the target, with no
gradient updates. \citet{zhou_context-driven_2025} attack the zero-history case
directly, retrieving the traffic of semantically similar pages to forecast a
new page that has no series of its own, which is the structural analogue of
forecasting a venue on its opening week from its siblings. GPHT
\citep{liu_generative_2024} shows the underlying mechanism, pretraining on a
mixed channel-independent corpus to reduce mean absolute error by about 6\% on
average relative to training from scratch.

Estate structure can also be exploited through reconciliation. MinT
\citep{wickramasuriya_optimal_2019} adjusts independent forecasts so that
venue-level predictions sum coherently to the estate total, minimising the
trace of the reconciled error covariance among all linear reconciliations that
preserve unbiasedness, an optimality that holds only if the base forecasts are
themselves unbiased --- a condition Section~\ref{sec:rw-ruler} shows is not
innocuous on a series of mostly zeros. HiGP \citep{cini_graph-based_2024} folds
the same coherence into an end-to-end model by projecting predictions onto the
space of coherent forecasts. The review by
\citet{athanasopoulos_forecast_2024} maps the resulting design space. These
results jointly justify treating a sparse or new venue not as an isolated
forecasting failure but as a borrower of rhythm from its estate and its
analogues.

\subsection*{Intermittent trade is a different object}

The literature above assumes demand arrives continuously. For a venue that
trades one or two days in a typical week it does not, and the distinction is
not cosmetic: it changes the model, the metric and the meaning of a good
forecast. \citet{croston_forecasting_1972} observed that exponential smoothing
applied to an intermittent series produces biased estimates, and proposed
overcoming the difficulty by making separate estimates of the size of demand
and of the demand frequency, so that the interval between events is modelled
rather than smoothed away. \citet{syntetos_accuracy_2005} showed that
Croston's estimator is nonetheless biased, because the expectation of a ratio
is not the ratio of expectations, and proposed the approximately unbiased
estimator that scales Croston's by $1 - \alpha/2$, where $\alpha$ is the
smoothing constant used for the inter-demand intervals.

The classification deciding when these methods apply is due to
\citet{syntetos_categorization_2005}, who partition series by average
inter-demand interval $p$ and the squared coefficient of variation of demand
sizes $v$ into smooth, erratic, intermittent and lumpy regions.
\citet{kostenko_note_2006} contribute two distinct things, and it is worth
keeping them apart because they answer different questions. They show the
published cutoffs contain arithmetic errors, correcting them to $p = 4/3$ and
$v = 0.5$ in place of 1.32 and 0.49; and they propose a simpler and more
accurate rule for choosing between the two estimators, dividing the quadrant
diagonally so that the bias-corrected estimator is used whenever
$v > 2 - \tfrac{3}{2}p$ and Croston's otherwise.

The consequence for a venue near the corrected cutoff is narrower than it first
appears, and runs the opposite way to the obvious reading. The correction moves
the \emph{label} a series receives, since $p = 4/3$ rather than 1.32 is the
boundary of the region in which Croston's method can still win. It does not
follow that the \emph{choice} of estimator is finely balanced there, because
the selection rule is a diagonal in the $(p, v)$ plane rather than a threshold
on $p$ alone: for $p > 4/3$ the quantity $2 - \tfrac{3}{2}p$ is negative and
the bias-corrected estimator is preferred at every $v$, while for $p$ just
below $4/3$ the switching value of $v$ is close to zero, so any series with
appreciable variation in its demand sizes sits far from the boundary and the
rule selects unambiguously. The honest statement is therefore that the
corrected constant changes how a venue is described, that the estimator choice
for a venue with non-negligible $v$ is not in doubt either way, and that a $p$
estimated from one year of a single venue's trade carries sampling uncertainty
that this chapter does not quantify and that the methodology should.

A second tradition models the same structure as two decisions rather than one.
\citet{cragg_statistical_1971} proposed the two-part specification in which a
binary model governs whether the outcome is positive and a separate model
governs the amount conditional on occurrence, and \citet{mullahy_specification_1986}
developed the unrestricted hurdle form, estimating the occurrence and magnitude
components by separate maximum likelihood. The attraction for a booking-led
venue is direct: whether the venue trades tomorrow is a different question,
answerable from different information, from how much it takes if it does. The
limitation is equally direct, and is a limitation of data rather than of
method --- a hurdle model is only as good as the occurrence signal available to
it, and where that signal lives outside the dataset the two-part structure
buys nothing.

\subsection*{Weather, and the temptation of exogenous data}

Weather is the exogenous covariate a hospitality forecaster reaches for first.
The relevant evidence is indirect, and this chapter states the inference it is
making rather than presenting it as settled. In load forecasting,
\citet{haben_short_2019} report that, in contrast to high-voltage load
forecasts, temperature was not an important factor in the accuracy of their
low-voltage forecasts: often it had little to no effect, and in many cases
including temperature was actually detrimental to accuracy. Their proposed
explanation is a mechanism rather than a coincidence --- temperature is
strongly correlated with seasonality, so a model can learn temperature as a
surrogate for the seasonal pattern it already encodes, and produce larger
errors than it would with no temperature at all. The general lesson is that a
covariate correlated with a seasonal signal the model already captures can cost
more than it contributes, and that this is likelier at low aggregation, where
the seasonal pattern is proportionally larger and the weather signal noisier.
\citet{hertel_explainable_2026} are consistent with this: on their load data,
of the total attribution assigned across the feature set of a
covariate-informed Chronos-2 forecast, roughly nine-tenths falls on the load
history, with temperature accounting for about 4\% and irradiance about 3\%.

The transfer of that lesson to hospitality is an assumption, and it is
contestable in a specific way. Electricity demand responds to temperature
through heating and cooling, which is exactly the channel that duplicates
seasonality; a public house responds to weather through outdoor trade and
footfall, a channel with no such duplication and a plausible claim to
independent information. The position this chapter takes is therefore weaker
than "weather will not help": it is that a single site is not a region, that
weather competes with a strong recurring weekly pattern at this level of
aggregation, and that a weather covariate must be shown to earn its place
rather than assumed to. This dissertation's own weather experiment returns a
null --- no detectable effect at the number of evaluation folds available ---
which is consistent with a small positive contribution, with none, and with a
small harmful one alike, and which therefore neither confirms nor refutes the
expectation set here.

\subsection*{From a point to a band}

A point forecast, however accurate, cannot say what would count as surprising.
The conformal-prediction line replaces the point with an interval that carries
a distribution-free coverage guarantee. In its split form, which is the
variant used here, a score function is calibrated on held-out data and an
empirical quantile of those scores sets the interval width.
\citet{angelopoulos_conformal_2023} state the resulting guarantee in two
directions at once: coverage is bounded below by the nominal level and above
by that level plus a term of order one over the calibration-set size. Both
bounds are on \emph{marginal} coverage, in expectation over the calibration
draw, and the upper bound additionally requires the conformity scores to be
almost surely distinct. That last condition is worth stating plainly because it
fails on precisely the data this work is concerned with: on a series of mostly
zeros, tied scores are the norm rather than the exception, so the upper bound
should not be applied to such a venue without qualification.

Two consequences follow from \citeauthor{angelopoulos_conformal_2023}'s bounds,
and they must be kept apart. Where the upper bound applies, a band whose
\emph{expected} coverage sits far above the nominal level is not a cautious
band but a mis-calibrated one. The corresponding claim about an observed run of
data is much weaker, because the guarantee is on coverage in expectation over
the calibration draw while realised coverage on a finite window is a random
variable around it: a band that happens to contain every observation in a short
test window is an unremarkable outcome of a small sample, and a shortfall
measured on such a window is equally unremarkable. Coverage reported without
the width that produced it, and without the number of points it was measured
on, conceals both facts at once --- a point Section~\ref{sec:rw-ruler} recovers
from a different direction, in \citeauthor{kolassa_evaluating_2016}'s argument
that calibration and sharpness must be scored together.

That guarantee holds under exchangeability, which trading data has no
particular reason to satisfy, and the online variants exist to address exactly
this. Adaptive conformal inference \citep{gibbs_adaptive_2021} adjusts the
nominal level in response to realised coverage. EnbPI \citep{xu_conformal_2021}
discards the exchangeability requirement altogether, guaranteeing approximately
valid marginal coverage under mixing conditions on the error process rather
than distributional assumptions on the data, and calibrating by updating past
residuals through a sliding window without refitting the underlying model; its
successor \citep{xu_sequential_2023} makes the residual model explicit, and
conformal PID control \citep{angelopoulos_conformal_2023-1} recasts the
adjustment as feedback control. The practical guidance is collected by
\citet[a 2025 preprint]{stocker_gentle_2025}. Adaptivity carries a price that
\citet{zaffran_adaptive_2022} make explicit: on exchangeable scores the
adaptive update degrades efficiency increasingly with its step size, so a
method that adapts to a shift that never arrives is worse than one that does
not adapt at all. \emph{Calibration}, in this sense, is the property that a
stated 90\% interval contains the truth about 90\% of the time, neither more
nor less. Expressing a venue's rhythm as a calibrated band, rather than a
single expected number, is the design choice this section ends on, because once
the rhythm is an interval the question shifts from \emph{what will happen} to
\emph{what would count as surprising}.

\section{Scoring the Rhythm: Error Measures and Model Selection}
\label{sec:rw-ruler}

A rhythm that cannot be scored honestly cannot be defended, and two distinct
questions have to be answered before any forecast in this dissertation means
anything: which error measure expresses what a venue actually cares about, and
by what procedure one model is declared better than another. Both are usually
treated as preliminaries. On intermittent hospitality data they are the
substance.

The scaled measures are the natural candidates on series of different
magnitudes, but \citet{hewamalage_forecast_2023} show that the choice is not
neutral: measures built on absolute errors optimise for the median, and on
intermittent series that property makes a constant zero look like the best
available prediction. The same paper notes the companion hazard for
benchmark-scaled measures, that a naive benchmark scoring exact zeros on zero
actuals deflates the denominator those measures divide by. Both failure modes
bite hardest on precisely the venue this estate most needs a metric for --- one
whose median trading day takes nothing. The M5 competition
\citep{makridakis_m5_2022} answered this by scoring on a squared scaled error,
aggregated across the hierarchy in a value-weighted form, across a retail
corpus reported as substantially intermittent: 77.3\% of its series have an
average inter-demand interval above $4/3$, a figure the competition report
gives against the corrected \citet{kostenko_note_2006} threshold rather than
the original one. That squared base errors optimise for the mean rather than
the median is established by \citet{hewamalage_forecast_2023} rather than by
the competition report. \citet{kolassa_why_2020} argues more generally that the
measure encodes an implicit choice of which functional of the predictive
distribution is wanted --- naming the mean, the median and the $(-1)$-median as
the summaries different measures reward --- so that selecting a measure is
selecting a target and not merely a scale.

Three further results converge on the same conclusion from different
directions. First, coherence. \citet{kolassa_we_2023} states the cost in his
title, though the incompatibility is narrower, and more useful, than the title
suggests. Because the median of a sum is not generally the sum of the medians,
point forecasts minimising MAE, or MASE, which is a scaled MAE, are usually not
coherent. The exception is the instructive part: coherence follows when the
measure is minimised by the conditional mean, because the expectation is
additive, and squared error is the canonical member of that class. The class is
in fact wider --- any Bregman divergence is minimised by the mean, so the
Poisson and Tweedie deviances share the property, a generalisation this chapter
notes rather than one \citeauthor{kolassa_we_2023} makes --- but the conclusion
is unaffected. The conflict is a property of absolute-error measures, not of
error minimisation as such, and a mean-eliciting measure dissolves it rather
than trading against it. For an estate that reconciles venue forecasts to an
estate total, this is not an aesthetic point: it is the condition under which
the reconciliation of Section~\ref{sec:rw-rhythm} is coherent with the
objective the forecaster was fitted to.

Second, the limit of error as an instrument at all.
\citet{chatfield_all-zero_2007} report, from a full factorial study varying
lumpiness, demand variation and the ordering, holding and shortage costs, that
the lowest forecasting error does not necessarily lead to the lowest system
cost, and --- the sharper half of the finding --- that all-zero forecasts yield
the lowest cost when lumpiness is high, and are also best at mid-lumpiness if
the shortage cost is much higher than the holding cost. Where a series is lumpy
enough, error minimisation and cost minimisation come apart, and the honest
response is to change the objective rather than to keep tuning against a
measure that has stopped tracking what matters.

Third, the point forecast may be the wrong object.
\citet{kolassa_evaluating_2016} argues that the emphasis on point forecasts
misses the mark for count and intermittent demand, and that one should instead
characterise the entire predictive distribution, from which point, interval or
quantile forecasts can be extracted as required. His instrument is the proper
scoring rule, which assesses calibration and sharpness together rather than
either alone. This closes a loop opened at the end of
Section~\ref{sec:rw-rhythm}: a band reported by its coverage alone is scored on
calibration with sharpness suppressed. Coverage and width are one measurement,
not two, and a proper scoring rule is the literature's name for treating them
that way.

The four strands above make a case for a mean-eliciting, squared-scaled measure
on this estate. It should be said here, rather than left for a reader to
discover, that this dissertation's reported headline accuracy figures are
mean-absolute-scaled rather than squared-scaled, with an unscaled absolute
error used at the venue where no scale basis proved defensible. The argument
assembled in this section therefore runs against the measure the results
chapter reports, and the tension is a limitation of this work rather than a
point in its favour: the reasons are ones of comparability with a pipeline and
a set of frozen artefacts established before this argument was assembled, and
the discussion, not this chapter, is where that trade should be defended.
Nothing in the sections that follow depends on the headline measure having been
squared-scaled.

The second question is comparison. Declaring one model better than another on a
difference of means is not a finding, and the forecasting literature has said
so for thirty years. \citet{diebold_comparing_1995} formalised the test of
equal predictive accuracy between two forecast sequences;
\citet{harvey_testing_1997} showed the statistic is seriously over-sized in
moderate samples and supplied the finite-sample correction
$\left[\frac{n + 1 - 2h + n^{-1}h(h-1)}{n}\right]^{1/2}$, whose behaviour is
itself informative. The bracketed quantity is exactly zero at $h = n$ and
$h = n+1$, and negative for real $h$ between them, so for an evaluation window
as short as the horizon the correction factor vanishes and the corrected
statistic is identically zero: computable, but carrying no information
whatever. The underlying test is already strained there for a second reason,
since an $h$-step Diebold--Mariano variance estimator needs enough observations
to form $h-1$ autocovariances. Both are facts about the design rather than
defects of the correction, and this dissertation states where they apply rather
than working around them. \citet{hansen_model_2011} generalise from the
pairwise case to the set: a model confidence set contains the set of superior
models with asymptotic probability at least $1 - \alpha$, and, in their own
formulation, acknowledges the limitations of the data, such that uninformative
data yield a set with many models whereas informative data yield a set with
only a few. A retained set that contains nearly every candidate should
therefore be read not as a wide-but-informative result but as a procedure that
has failed to eliminate anything, which is a statement about the evidence
available rather than a finding of equivalence.

Two recent studies show how fragile rankings are when this discipline is
absent. \citet{brigato_there_2025}, evaluating eight models on fourteen
datasets across roughly five thousand trained networks, find that slight
changes to experimental setups or evaluation metrics drastically shift the
common belief that newly published results advance the state of the art, and
that even minor deviations in setup can dramatically change performance
rankings. \citet{hewamalage_look_2021} make the same concern measurable,
proposing rank stability --- how much the ranking of an experiment changes
across similar datasets when models and errors are held constant --- and
finding that the main drivers of instability are hierarchical aggregation and
scaling, with the scale normalisation of the M5 error measure less stable than
other scale-free errors. Three things follow. The number of evaluation origins
is not a free parameter, because a ranking computed on a handful of origins
carries little information about the ranking on hundreds. The case for a
squared scaled error assembled above is strong but not costless, since the
scale normalisation it depends on is itself a documented source of rank
instability. And --- the limit of what this literature licenses --- these
results predict that a ranking may move, not which way it will move. This is
the first of the two places where the literature and this project's
observations point in different directions: the reversal this dissertation
reports on increasing its evaluation origins moved in the direction that
favoured the model already in production, and no work cited here predicts that
direction.

\section{Noticing Deviation: From Calibrated Forecast to Signal}
\label{sec:rw-deviation}

Once the rhythm is a calibrated interval rather than a point, a meaningful
\emph{deviation} is definitionally a breach of that band, and the detection
literature's job is to make breach-detection reliable and honestly evaluated
on short, noisy series. The classical methods organise by what they detect.
Bayesian online change-point detection \citep{adams_bayesian_2007} maintains a
posterior over the run length since the last change, suited to abrupt regime
shifts such as a venue's permanent closure; the offline survey by
\citet{truong_selective_2020} reduces the whole family to a choice of cost
function, search method and penalty, which clarifies that change-point
detection answers \emph{when the regime moved}. That reduction also locates
the cumulative-sum scheme of \citet{page_continuous_1954}, the sequential
counterpart to the offline cost functions the survey enumerates. Page assigns
a score to each observation, accumulates those scores, and signals when the
running total has risen a fixed amount above its own running minimum, which
makes the scheme sensitive to small sustained shifts that a per-observation
test will not register. He also supplies the criterion by which such a
detector should be judged. The average run length, which he defines as the
expected number of observations before action is taken, measures the expense
of false alarms when conditions are stable and the delay before a real change
is acted on when they are not. That two-sided cost is the same one this
chapter arrives at independently in Section~\ref{sec:rw-evaluation}, stated
for an industrial inspection scheme seventy years before it was restated for
conversational agents. Being cheap to evaluate on each new observation and
requiring no posterior over run length, the cumulative-sum scheme is the
detector this work puts into production, with Bayesian online change-point
detection retained as the offline benchmark it is scored against.
Point-anomaly methods answer a different question, \emph{is this observation
extreme}, and SPOT \citep{siffer_anomaly_2017} answers it by fitting a
generalised Pareto tail via extreme value theory, requiring only a risk level
and assuming nothing about the data distribution. Mapping these onto a venue is
straightforward: a sudden quiet Saturday is a point anomaly against the band, a
venue closure is a change point, and conflating the two is a modelling error.

The harder problem is evaluation, and here the literature contains a finding
that should unsettle any reported detection score. \citet{kim_towards_2022}
show that the point-adjustment protocol, which counts a whole anomalous segment
as detected if any single point inside it crosses the threshold, is so generous
that randomly generated anomaly scores reach an adjusted F1 near one and
overturn most state-of-the-art results. \citet{liu_elephant_2024} confirm and
extend this with TSB-AD, a curated benchmark of 1070 series, identifying VUS-PR
as the measure that stays stable under small detection lags while remaining
unbiased against random scores, and finding along the way that simpler
statistical detectors often beat elaborate neural ones, which echoes the
baseline lesson of Section~\ref{sec:rw-rhythm}. Corrective protocols exist that
penalise the false positives the standard adjustment ignores, whether by
balancing the adjustment or by decaying credit for late detection
\citep{bhattacharya_towards_2024,gim_evaluation_2023}. The implication for this
work is methodological rather than incidental: detection on hospitality data
should be reported with a lag-tolerant, random-robust measure in the spirit of
VUS-PR, and point-adjusted F1 is rejected as a headline metric because it
would flatter a detector that has learned nothing.

The thread of Section~\ref{sec:rw-rhythm} closes here. CPTC
\citep{sun_conformal_2025} fuses the two halves directly, pairing a switching
dynamical model that predicts the underlying state with online conformal
prediction calibrated separately per state. Its guarantees repay precise
statement, because they are narrower than the framing invites. Exact
finite-sample coverage holds only under exchangeability, which a change point
violates by construction. What survives a regime shift is time-averaged
asymptotic validity, resting on an assumption that the states themselves have
a stationary distribution, together with faster post-shift convergence than
adaptive conformal inference, which the authors describe as shortening the
miscoverage period rather than removing it. The residual risk then concentrates
in a single quantity, since coverage degrades in proportion to the rate at
which the state is misclassified. \citet{barber_conformal_2023} give the
general form of the same trade, bounding the loss of coverage under arbitrary
departures from exchangeability by a weighted sum of total-variation distances
between residual vectors, and doing so with no assumption whatever on the joint
distribution of the data.

Two qualifications belong here, because both cut against the arc this section
has built. The first is that the case for online adaptation is theoretical and
conditional, not empirical and general. The guarantees above concern rates of
convergence after a shift; none of them promises that an adaptive procedure
will out-perform a static one on any particular series, and
\citet{zaffran_adaptive_2022}'s efficiency penalty runs the other way. This is
the second of the two places where the literature and this project's
observations diverge: at the one genuine regime change in this estate's data,
adaptive calibration performed \emph{worse} than leaving the band alone, which
the framing of this section would not have led a reader to expect. The second
qualification concerns under-coverage. A split-conformal band on this estate's
anchor venue covers below its nominal level. It is tempting to read that as
evidence of an exchangeability violation, and under the right conditions it
would be, but the inference does not go through on a finite window: the
guarantee is on coverage in expectation over the calibration draw, and realised
coverage varies around that expectation, falling below nominal a substantial
fraction of the time even when the assumptions hold exactly. What such a
shortfall means therefore depends on the width of the band, the number of
points it was measured on, and the sampling interval around the coverage
estimate --- which is the reporting standard this chapter has already argued
for on other grounds.

Read together, these results reposition the design question this chapter has
been building toward. A band cannot be promised to hold through a regime
change. It can only be promised to lose coverage in proportion to how badly
the regime is estimated. For a hospitality estate the operational consequence
is to prefer a regime variable that is observed over one that is inferred,
because whether a venue is trading is recorded in the till rather than hidden
in the residuals, and an observed state drives the misclassification term to
zero. A detector can establish \emph{that} trade has broken from its rhythm;
it cannot decide whether the break is worth a manager's attention.

\section{Reasoning and Surfacing: Agents that Act Without Being Asked}
\label{sec:rw-surfacing}

Converting a deviation signal into a timely, well-judged intervention needs
more than tool-use plumbing; it needs an agent that initiates action and
decides what is worth raising. The plumbing itself is mature: interleaved
reasoning and tool calls, self-supervised API invocation, and verbal
self-correction across attempts are established primitives
\citep{yao_react_2022,schick_toolformer_2023,shinn_reflexion_2023}, they are
present in the existing product, and they are not where the open problem lies.
That the problem is open is visible in evaluation: $\tau$-bench
\citep[a 2024 preprint]{yao_-bench_2024} places agents in multi-turn
tool-and-user interactions and finds frontier models far from solving them,
with even GPT-4o completing only 35.2\% of airline-domain tasks on a
single-attempt basis, so reliability rather than capability is the binding
constraint even before proactivity enters.

Memory is the bridge from the forecasting half of this work to the agent half,
and the connection must be stated carefully, because it is easy to claim more
for it than is warranted. Retrieval-augmented generation
\citep{lewis_retrieval-augmented_2020} established the template of grounding a
generator in an external, non-parametric store accessed by a learned
retriever. \citet[a 2026 preprint]{hu_memory_2026} argue that agent memory is
categorically more than this: where classical retrieval reads from a static
corpus per query, agent memory is a persistent, self-evolving cognitive state
that accumulates the agent's own experience over time. Generative agents
\citep{park_generative_2023} make the mechanism concrete, with a memory stream
surfaced by a weighted combination of recency, relevance and importance, and
recursive reflection --- triggered when accumulated importance crosses a
threshold, rather than on a clock --- that abstracts raw records into
higher-level inferences.

The reframing this work takes from that literature is deliberately bounded. A
learned rhythm is a durable, evolving record of normality that the agent
retrieves against when deciding whether the present is worth remarking on, and
in that specific respect it does the work of the retrieval half of a memory
stream. What a forecasting band does not supply is the rest of the
architecture: there is no reflection step abstracting episodes into
higher-level inferences, and no store of the agent's own past interventions to
retrieve from. The claim made here is accordingly the narrow one --- the rhythm
serves as the agent's model of normality, not as a full memory system in the
sense \citet{hu_memory_2026} and \citet{park_generative_2023} describe.
Retrieval over live operational data carries a caveat in either case, since
such stores are an attack surface, as PoisonedRAG
\citep{zou_poisonedrag_2025} demonstrates, which matters once the memory is
populated from production traffic.

The frontier of proactive agency is where this work positions its contribution,
and the evidence base here is the least settled in the chapter: every result in
the next two paragraphs is a preprint or an unrefereed submission.
\citet[a 2024 preprint]{lu_proactive_2024} provide the template. Their
ProactiveBench gathers 6,790 training and 233 test events from real human
activity, trains an agent to anticipate needs, and, tellingly, the dominant
failure they document is over-offering. Their false-alarm rate is the
proportion of predictions that were not needed, and on that definition every
strong proprietary model they evaluate exceeds one half, from 51.85\% for
GPT-4o to 64.73\% for GPT-4o-mini. Their best fine-tuned agent reaches an F1 of
66.47\% and a separately trained reward model attains 91.80\% agreement with
human judges; on a 233-event test set the sampling uncertainty on an agreement
rate of that size is around two percentage points, so the agreement figure
should be read as approximately 92\% rather than to the precision stated, and
the F1, whose sampling distribution is not that of a simple proportion, should
be read as approximately 66\%. The headline finding is not that proactivity is
hard to produce but that it is hard to \emph{restrain}.

Subsequent work attacks restraint from different angles. ProActor
\citep[an unrefereed conference submission]{ding_proactor_2026} reframes timing
as a window of valid moments rather than a single labelled point and optimises
it with turn-level reinforcement learning, on the argument that penalising any
deviation from one rigid timestamp is the wrong objective. ProAgentBench
\citep[a 2026 preprint]{tang_proagentbench_2026} grounds the question in 28,000
events from over 500 hours of real sessions and states the cost structure
plainly: low precision in deciding when to assist causes alert fatigue, where
excess false alarms drive users to ignore or disable the feature. PRISM
\citep[a 2026 preprint]{fu_prism_2026} draws the consequence, formulating
intervention as a cost-sensitive decision in which the agent speaks only when
calibrated acceptance probability clears a threshold set by the asymmetric
costs of false alarms and misses. Measured against
\citeauthor{lu_proactive_2024}'s best fine-tuned agent as its baseline, it
raises F1 from 66.47\% to 86.61\% --- a gain of about twenty percentage points
--- while reducing the false-alarm rate from 50.22\% to 22.94\%, a relative cut
of roughly one half. Both movements follow from the same gating threshold and
should be read as one result rather than two.
\citet[a 2026 preprint]{gulati_ask_2026} add that the urgency of a clarifying
question depends on what is missing: goal clarification decays quickly,
retaining near-oracle value only if asked very early, while input clarification
remains recoverable through roughly half of execution, so a single uniform
urgency for every kind of missing information is wrong. Adjacent systems for
dialogue, wearables and mobile interfaces
\citep{liu_proactiveeval_2025,yang_contextagent_2025,yang_fingertip_2025}
confirm proactivity is being pursued across domains without supplying a
precedent in a multi-venue operational setting. Read together, this body of
work converges on one diagnosis that the next section takes up: the binding
failure of a proactive agent is the false alarm, and the open question is how
to weigh it.

\section{Judging the Judgement: Evaluating Proactive Intervention}
\label{sec:rw-evaluation}

The agent's decision to speak is only as good as our ability to tell a useful
prompt from an unwelcome one, and accuracy is the wrong instrument for that
measurement. Two literatures supply the right one. The first concerns
automated evaluation. \citet{zheng_judging_2023} show that a strong language
model as judge agrees with human preference above 80\%; on MT-bench pairwise
comparisons with ties excluded, GPT-4 reaches 85\% agreement with humans while
humans agree with each other 81\% of the time. The two rates are computed over
different pairings and are not a paired comparison, so the defensible reading is
that the judge performs about as well as a human annotator rather than better.
The qualifications are serious in any case. \citet{wang_large_2024} demonstrate
a positional bias severe enough that swapping the order of two candidates can
reverse the verdict; \citet{panickssery_llm_2024} show that judges recognise
and favour their own generations, with self-preference scaling with
self-recognition; and JUDGE-BENCH \citep{bavaresco_llms_2025} finds substantial
variance in agreement across models and datasets, with agreement depending on
the property being evaluated, the expertise of the human judges, and whether
the text being judged is human- or model-generated, concluding that judges must
be carefully validated against human judgments before they are used as
evaluators. The position this work takes is to use an automated judge only as
an instrumented, bias-audited proxy, with the real target being human
adopt-or-dismiss outcomes on actual surfaced insights.

The second literature defines what the metric must encode, and it is the
human-factors account of trust in automation. \citet{meyer_conceptual_2004}
draws the distinction that sets the loss function: \emph{compliance} is the
operator's response to a warning, \emph{reliance} the response to its absence,
and crucially the two are governed by different properties of the detector.
Compliance depends on the response criterion, responses to warnings becoming
less likely when the probability of an unjustified warning is high, the
cry-wolf syndrome in which operators cease or slow their response. Reliance, on
his account, depends instead on the system's sensitivity.

That the two failures are not merely different but unequal is a further claim,
and it has been tested directly. \citet{dixon_independence_2007} placed
operators under diagnostic automation that was either false-alarm-prone or
miss-prone and found that false-alarm-prone automation hurt overall performance
more than miss-prone automation, and that it degraded both compliance and
reliance while miss-prone automation appeared to affect only reliance. Their
stated implication for design is the premise this dissertation's loss function
encodes: false alarms are more damaging to overall performance than misses. The
evidence is a laboratory study of thirty-two participants on a dual tracking
and monitoring task, reported as a direction without an effect size, and the
asymmetry is asserted here at that strength and no higher; the authors
themselves conclude only that compliance and reliance are not entirely
independent.

Field evidence puts a magnitude on the same mechanism.
\citet{ancker_effects_2017}, studying alerts issued by 112 ambulatory clinicians
over three and a half years of electronic health record data, estimated that the
\emph{odds} of a clinician accepting a reminder fell by about 30\% for each
additional reminder received in the same encounter, and by about 10\% for each
five-percentage-point increase in the proportion of reminders that were
repeats. These are odds ratios from a multivariable model fitted over the full
set of alerts rather than changes in acceptance probability, and the
corresponding change in probability depends on the baseline acceptance rate.
The setting is clinical rather than hospitality and the numbers do not transfer,
but the mechanism does, and it is the mechanism a proactive agent must budget
against: the cost of an unnecessary alert is not paid by that alert, it is paid
by the next one. The structure is not new, and the detection literature reached
it first --- Page's average run length \citep{page_continuous_1954} already
traded the expense of false alarms under stable conditions against the delay
before a real change is acted on. What the human-factors work adds is that in a
managerial setting the two costs are borne by a person whose willingness to
keep reading is itself consumed by the false alarms.

The remaining literature situates that asymmetry in a theory of appropriate
use. \citet{lee_trust_2004} frame the goal as appropriate reliance, calibrated
trust in which trust matches capability, with overtrust producing misuse and
distrust producing disuse, the inappropriate rejection of a capable aid that
\citet{parasuraman_humans_1997} tie directly to high false-alarm rates.
\citet{parasuraman_complacency_2010} document the opposite pole, automation
bias and complacency. In the correlational half of their meta-analysis,
\citet{hancock_meta-analysis_2011} find robot performance-based factors the
stronger driver of trust, at $\bar{r} = +0.34$ with a confidence interval of
$+0.25$ to $+0.43$ against $\bar{r} = +0.03$ with an interval of $-0.09$ to
$+0.15$ for attribute-based factors; the intervals do not overlap, so the
ordering is established for that analysis. Their experimental half is weaker
and points less cleanly, resting the performance estimate on two studies
against eight for attributes, which is why the ordering is quoted here from the
correlational analysis alone.

The consequence for evaluation is that the right metric is calibrated precision
against real dismissals, weighted by the asymmetric cost of an interruption,
not an undifferentiated F1. A near-ready template exists: HiL-Bench
\citep[a 2026 preprint]{trinh_hil-bench_2026} scores agents with Ask-F1, the
harmonic mean of question precision and blocker recall, which penalises both
over-asking and silent wrong guessing; adapted from its coding domain to
adopt-versus-dismiss outcomes, it would capture the two-sided failure a
hospitality agent must avoid. Two conditions govern whether such an adaptation
measures anything. A two-sided metric is informative only where the evaluation
contains both failure modes, since a sweep over cost ratios in which one side
never occurs is degenerate whatever value it returns. And the calibration half
must actually be measured: \citet{guo_calibration_2017} show modern networks
are overconfident and that temperature scaling restores calibration as measured
by expected calibration error, which is the natural instrument for an agent's
stated confidence and the one that would let the coverage guarantee of
Section~\ref{sec:rw-deviation} and the calibration-with-sharpness argument of
Section~\ref{sec:rw-ruler} reappear as an evaluation guarantee on the agent's
output. That instrument is identified by this chapter and is not computed in
this dissertation, which is a limitation of the work and is recorded as one.

\section{Synthesis: What the Literature Leaves Open}
\label{sec:rw-synthesis}

Four results stand established by the work above, and each with a qualification
this chapter has argued for rather than inherited. Operational rhythm can be
borrowed across venues and analogues despite scarce data
\citep{das_-context_2025,zhou_context-driven_2025,liu_generative_2024}, though
every such demonstration is on a collection far larger than a three-venue
estate, and the theory of why pooling pays is a theory about large collections
\citep{montero-manso_principles_2021}. A deviation can be defined as the breach
of a calibrated band whose loss of coverage under a regime change is bounded
rather than eliminated
\citep{sun_conformal_2025,barber_conformal_2023,liu_elephant_2024}. An agent
can be made to act unprompted, with its learned rhythm serving as the model of
normality it retrieves against, though not as a full memory architecture
\citep{lu_proactive_2024,hu_memory_2026,park_generative_2023}. And the worth of
such an action is measurable against the asymmetric cost of a false alarm, an
asymmetry that is empirically tested rather than merely assumed
\citep{meyer_conceptual_2004,dixon_independence_2007,ancker_effects_2017,fu_prism_2026,trinh_hil-bench_2026}.
Each has been demonstrated in isolation, and in domains --- coding, dialogue,
web traffic, wearables, clinical informatics --- distant from hospitality
operations.

The fourth demands the most care, because the method it calls for already
exists. PRISM \citep{fu_prism_2026} sets its intervention threshold from an
asymmetric miss-to-false-alarm cost ratio, and among the proactive systems
surveyed here it is the one that gates on a calibrated acceptance probability
rather than on accuracy alone. It reports the gate's effect but not the
calibration of the probability the gate depends on. No claim of methodological
novelty in cost-sensitive intervention is therefore available to this
dissertation, and none is made. What PRISM leaves undone, and what none of the
surveyed systems attempts, is to fix that cost ratio to the stated preferences
of the person who bears the cost, to measure the calibration of the acceptance
probability that ratio is applied to, and then to score the resulting decisions
against that same person's accept-or-dismiss judgements on signals drawn from
their own business. The proactive systems surveyed in
Section~\ref{sec:rw-surfacing} are evaluated instead against annotated corpora
of recorded activity \citep{lu_proactive_2024,tang_proagentbench_2026}, against
simulated or scripted users
\citep{yao_-bench_2024,liu_proactiveeval_2025,gulati_ask_2026}, or against a
language model acting as judge \citep{fu_prism_2026,ding_proactor_2026}; none
is scored against the decisions of the operator whose work it is intervening
in. It should be said plainly that this positioning rests on a 2026 preprint: if
PRISM's reported results do not survive review, the gap this dissertation
occupies is differently shaped, though the absence of an operator-grounded
evaluation across the systems just enumerated would remain.

The contribution is accordingly one of field instantiation rather than of
method. \citet{paleyes_challenges_2022}'s survey is assembled from reports of
systems in deployment and finds practitioners encountering difficulties at
every stage of the deployment workflow --- a body of knowledge that exists
because deployments were written up, not because benchmarks were run. The
contribution attempted here is to take a risk-sensitive intervention policy of
the kind PRISM specifies, place it above a rhythm borrowed across a small real
estate and a deviation signal defined by a calibrated band, and evaluate the
assembly against an operator's own decisions in a live multi-venue setting.

The scope of any such claim is narrow and is stated here rather than deferred.
The setting is a single estate under a single operator, so the accept-or-dismiss
judgements available come from one rater, with no inter-rater reliability. The
estate comprises one anchor venue with substantial trading history, one sparse
booking-led venue whose occurrence signal is absent from the dataset, and one
venue that closed part-way through the study period, imposing a structural
break; a fourth location, immaterial by value, is excluded from all modelling.
The record covers a single seasonal cycle, so no year-over-year seasonality is
estimable at all. Findings from a design of this shape are existence proofs and
failure modes, not estimates that generalise to hospitality operations.

What this chapter has additionally tried to do is state where the literature
leads and where it does not. The locality--globality principle, extended by
analogy, suggests that pooling across three venues will buy little
\citep{montero-manso_principles_2021}; the low-voltage load-forecasting evidence
suggests that a weather covariate must earn its place rather than be assumed to
\citep{haben_short_2019,hertel_explainable_2026}; and the rank-stability
literature warns that a ranking computed on a handful of evaluation origins
carries little information, without predicting which way it will move
\citep{hewamalage_look_2021,brigato_there_2025}. Against those, two of this
project's observations run the other way from the chapter's own reading of the
literature, and both are marked where they arise: the direction of the ranking
reversal in Section~\ref{sec:rw-ruler}, and the performance of adaptive
calibration at a regime change in Section~\ref{sec:rw-deviation}. Whether a
policy developed on annotated benchmarks survives contact with one operator,
three venues and a year of untidy till data is the question this dissertation
sets out to answer, within the limits just stated.
